\documentclass[11pt,a4paper]{article}
\usepackage[T1]{fontenc}
\usepackage[utf8]{inputenc}
\usepackage[french]{babel}
\usepackage{geometry}
\geometry{margin=2.5cm}
\usepackage{xcolor}
\usepackage{booktabs}
\usepackage{array}
\usepackage{tikz}
\usetikzlibrary{shapes.geometric,arrows.meta,positioning,fit}
\usepackage{fancyhdr}
\usepackage{hyperref}
\usepackage{listings}
\usepackage{amsmath}

\hypersetup{colorlinks=true,linkcolor=blue,urlcolor=blue,citecolor=blue}

\lstset{
  basicstyle=\ttfamily\small,
  breaklines=true,
  frame=single,
  backgroundcolor=\color{gray!8},
  columns=fullflexible,
}

\setlength{\headheight}{14pt}
\pagestyle{fancy}
\fancyhf{}
\fancyhead[L]{Descripteur de percolation ICOHP/ICOBI}
\fancyhead[R]{\thepage}
\renewcommand{\headrulewidth}{0.4pt}

\title{\bfseries Descripteur de percolation de moindre résistance \\
à partir des données ICOHP/ICOBI \\
\large Objectif, méthodologie, workflow et environnement de calcul}
\author{Gilles Frapper \\ Professeur des Universités, Chimie Théorique \\ IC2MP UMR 7285 CNRS, Université de Poitiers \\ \href{mailto:gilles.frapper@univ-poitiers.fr}{gilles.frapper@univ-poitiers.fr}\\[4pt] \small Rapport généré dans le cadre du projet \texttt{viability}}
\date{13 août 2026}

\begin{document}
\maketitle

\tableofcontents
\newpage

\section{Objectif et problématique}

L'objectif est de disposer d'un \textbf{descripteur statique, calculable en
post-traitement}, qui approxime la voie de rupture structurale la plus
probable d'un cristal, sans recourir à une dynamique moléculaire ni à un
calcul NEB (nudged elastic band), et sans artefact lié à la taille d'une
supercellule.

Le point de départ est une structure cristalline relaxée (maille primitive)
et l'analyse de population de recouvrement associée (ICOHP et/ou ICOBI,
issue de LOBSTER), qui fournit pour chaque paire d'atomes en interaction :
la force de liaison, la distance, et le \textbf{vecteur de translation de
réseau} reliant les deux atomes (c'est-à-dire si la liaison reste dans la
maille de référence ou relie un atome à l'image périodique d'un voisin dans
une maille adjacente).

L'intuition physique est la suivante : la voie de moindre résistance à
travers le cristal, dans une direction cristallographique donnée, est la
séquence de liaisons dont la somme des forces traversées est minimale, tout
en formant un chemin qui \emph{traverse effectivement} la maille (et non un
simple aller-retour local). Un tel chemin, en théorie des graphes
périodiques, est appelé \textbf{cycle non contractile} : c'est une marche
dans le graphe qui relie un atome à l'une de ses propres images
périodiques, avec une somme des vecteurs de translation traversés
exactement égale à un vecteur de réseau primitif.

Ce problème est mathématiquement identique à celui rencontré dans
l'analyse de percolation ionique des électrolytes solides (recherche de
chemins de migration à travers un réseau cristallin), mais appliqué ici à
la force de liaison (ICOHP/ICOBI) plutôt qu'à une barrière de migration.

\section{Méthodologie}

\subsection{Graphe périodique étiqueté (voltage graph)}

Contrairement à une approche par supercellule (qui duplique physiquement
les atomes et est sujette à des artefacts de taille finie), la périodicité
est ici gérée par un \textbf{étiquetage vectoriel des arêtes} :

\begin{itemize}
  \item les n\oe uds du graphe sont les atomes de la maille primitive
    uniquement (aucune duplication géométrique) ;
  \item chaque liaison ICOHP/ICOBI devient une arête étiquetée par le
    vecteur de translation entier $(n_x, n_y, n_z)$ qu'elle traverse ---
    $(0,0,0)$ si les deux atomes sont dans la même maille, $(1,0,0)$ si la
    liaison relie un atome à l'image de son voisin translatée d'un vecteur
    de réseau selon $\mathbf{a}$, etc. ;
  \item le poids de l'arête est $|\text{ICOHP}|$ (ou $|\text{ICOBI}|$) ---
    sans inversion : on cherche le chemin dont la somme des forces de
    liaison traversées est \emph{minimale}.
\end{itemize}

\subsection{Recherche du cycle non contractile de poids minimal}

Pour chacun des trois vecteurs de réseau primitifs $\mathbf{a}$,
$\mathbf{b}$, $\mathbf{c}$, on cherche le chemin de poids minimal dans
l'espace d'état étendu $(\text{atome}, \text{translation cumulée})$ reliant
l'état $(\text{atome}, (0,0,0))$ à l'état $(\text{atome}, \text{direction})$,
pour chaque atome de départ possible de la maille. Les poids d'arête étant
positifs ou nuls (valeurs absolues), c'est un problème de plus court chemin
classique, résolu par l'algorithme de \textbf{Dijkstra} sur ce graphe
d'état étendu :

\begin{lstlisting}
etat = (atome, (tx, ty, tz))
depart = (atome, (0, 0, 0))
arrivee = (atome, direction)   # direction in {a, b, c}
# Dijkstra standard sur le graphe des etats,
# translation cumulee bornee par --coord-bound
\end{lstlisting}

Le minimum sur les atomes de départ donne le poids de percolation de la
direction considérée ; le minimum sur les trois directions donne le
descripteur principal. L'espace des translations cumulées est borné
(paramètre \texttt{--coord-bound}) ; cette borne est dérivée automatiquement
du plus grand vecteur de translation présent dans les données d'entrée ---
elle est donc directement liée au rayon de coupure du calcul ICOHP/ICOBI en
amont, et n'est \emph{pas} un paramètre de convergence à optimiser.

\subsection{Cas déconnecté}

Si aucun état-cible n'est atteint dans une direction donnée (réseau de
liaisons déconnecté selon cette direction, compte tenu du rayon de coupure
ICOHP/ICOBI), le résultat est explicitement marqué \texttt{disconnected},
jamais silencieusement remplacé par une valeur infinie ou nulle.

\subsection{Sortie}

Pour chaque composé : le poids de percolation minimal toutes directions
confondues (descripteur principal), la direction correspondante, les poids
par direction $a/b/c$ avec leur statut, ainsi que les agrégats classiques
(somme, moyenne, min, max des ICOHP/ICOBI) pour comparaison directe.

\section{Organigramme du workflow}

\begin{center}
\begin{tikzpicture}[
  node distance=0.9cm and 1.6cm,
  box/.style={draw, rounded corners, minimum width=4.6cm, minimum height=1.05cm,
              align=center, fill=blue!6, font=\small},
  io/.style={draw, minimum width=4.6cm, minimum height=1.05cm, align=center,
             fill=orange!10, font=\small},
  arr/.style={-{Latex[length=2.2mm]}, thick}
]
\node[io]  (mp)     {Materials Project (API) \\ structures PBE, exp\'erimentales};
\node[box, below=of mp] (fetch)   {\texttt{fetch\_candidates.py} \\ tri hull / m\'etastable $\times$ ionique / covalent / m\'etallique};
\node[box, below=of fetch] (dl)     {\texttt{download\_selected.py} \\ POSCAR relax\'e (pymatgen)};
\node[box, below=of dl] (prep)   {\texttt{prepare\_vasp\_lobster.py} \\ INCAR / KPOINTS / POTCAR / lobsterin};
\node[box, below=of prep] (vasp)   {VASP (statique, ISYM=$-1$) \\ WAVECAR, CHGCAR};
\node[box, below=of vasp] (lob)    {LOBSTER 5.1.1 \\ ICOHPLIST.lobster, ICOBILIST.lobster};
\node[box, below=of lob] (perc)   {\texttt{percolation\_path.py} \\ graphe p\'eriodique \'etiquet\'e + Dijkstra};
\node[io, below=of perc] (out)    {\texttt{results.csv} / \texttt{results.json} \\ par composé};

\draw[arr] (mp) -- (fetch);
\draw[arr] (fetch) -- (dl);
\draw[arr] (dl) -- (prep);
\draw[arr] (prep) -- (vasp);
\draw[arr] (vasp) -- (lob);
\draw[arr] (lob) -- (perc);
\draw[arr] (perc) -- (out);
\end{tikzpicture}
\end{center}

\section{Environnement et installation}

\subsection{Machine de calcul}

Cluster \textbf{Yargla} (IC2MP, Université de Poitiers,
\texttt{yargla.chimie.univ-poitiers.fr}), ordonnanceur SLURM, partition
\texttt{defq}.

\subsection{Logiciels}

\begin{itemize}
  \item \textbf{VASP} 6.5.0 (module \texttt{vasp/6.5.0}), potentiels
    PAW--PBE recommandés par élément.
  \item \textbf{LOBSTER} 5.1.1 (module \texttt{lobster/5.1.1}) --- version la
    plus récente disponible sur le cluster.
  \item \textbf{Python} 3.11, environnement virtuel dédié
    \texttt{/home/gilles/viability/.venv} :
    \texttt{pymatgen} (analyse structurale, parsing LOBSTER),
    \texttt{lobsterpy} (utilitaires LOBSTER),
    \texttt{mp-api} (requêtes Materials Project).
\end{itemize}

\subsection{Installation (résumé)}

\begin{lstlisting}
module load intel
module load impi/2021.13
module load vasp/6.5.0
module load lobster/5.1.1

python3.11 -m venv .venv
source .venv/bin/activate
pip install pymatgen lobsterpy mp-api
\end{lstlisting}

Chaque calcul VASP+LOBSTER est soumis via un script SLURM
(\texttt{submit.sh}, 1 n\oe ud, 16 t\^aches, 24h de mur maximum), qui
enchaîne \texttt{vasp\_std} puis \texttt{lobster-5.1.1} dans le même
répertoire de travail.

\section{Jeu de données}

Six composés expérimentaux (référencés ICSD), relaxés PBE par Materials
Project, $\leq 20$ atomes/maille, répartis en 2 familles $\times$ 3 types
de liaison :

\begin{center}
\begin{tabular}{@{}lllr@{}}
\toprule
Famille & Type de liaison & Composé (mp-id) & $E_\text{hull}$ (eV/atome) \\
\midrule
sur le hull   & ionique    & NaCl (mp-22862)              & 0.0000 \\
sur le hull   & covalent   & Si (mp-149)                  & 0.0000 \\
sur le hull   & métallique & AlNi, B2 (mp-1487)           & 0.0000 \\
métastable    & ionique    & LiBr (mp-23259)               & 0.0099 \\
métastable    & covalent   & C, graphite rhomboédrique (mp-169) & 0.0008 \\
métastable    & métallique & BeCu (mp-2323)                & 0.0050 \\
\bottomrule
\end{tabular}
\end{center}

\section{Résultats}

\begin{center}
\small
\begin{tabular}{@{}lrrrr@{}}
\toprule
Composé & Poids percolation & Direction & $\text{ICOHP}_\text{min}$ & $\text{ICOHP}_\text{moy}$ \\
        & (eV) & & (eV) & (eV) \\
\midrule
NaCl (hull, ionique)      & 0.0307 & $a$ & $-0.595$  & $-0.089$ \\
Si (hull, covalent)       & 0.0032 & $c$ & $-4.552$  & $-0.291$ \\
AlNi (hull, métallique)   & 0.0042 & $a$ & $-1.992$  & $-0.322$ \\
LiBr (métastable, ionique)     & 0.0988 & $c$ & $-2.041$  & $-0.338$ \\
C rhomb. (métastable, covalent) & 0.0003 & $c$ & $-11.737$ & $-0.286$ \\
BeCu (métastable, métallique)   & 0.0005 & $b$ & $-1.848$ & $-0.162$ \\
\bottomrule
\end{tabular}
\end{center}

\section{Discussion : le descripteur diverge fortement des agrégats classiques}

Le résultat le plus marquant est que, pour les six composés, le poids de
percolation est \textbf{un à quatre ordres de grandeur plus faible} que
l'ICOHP minimal (liaison la plus forte) et même que l'ICOHP moyen. C'est
exactement le comportement recherché : le descripteur de percolation ne
mesure pas la force de la liaison la plus forte, mais la résistance du
chemin \emph{le plus faible} qui traverse effectivement le cristal --- et
ce chemin n'emprunte pas nécessairement les liaisons les plus fortes.

\textbf{Explication physique pour NaCl (cas type).} La liaison Na--Cl la
plus forte ($-0.595$~eV) relie les deux atomes \emph{à l'intérieur de la
même maille primitive} (vecteur de translation $(0,0,0)$) : dans la maille
primitive rhomboédrique fournie par Materials Project pour cette structure
de type rock-salt, cette liaison ne traverse donc \emph{aucune} frontière
de maille et ne peut à elle seule contribuer à un cycle non contractile.
Le chemin non contractile le moins coûteux emprunte donc des interactions
Na--Na ou Cl--Cl de deuxième sphère de coordination ($d \approx 3.95$~\AA,
$|\text{ICOHP}| \approx 0.031$~eV), nettement plus faibles mais qui, elles,
traversent effectivement une maille dans la direction considérée.

\textbf{Implication méthodologique.} Ce comportement est correct au sens
strict de la définition topologique demandée (cycle non contractile de
poids minimal sur les vecteurs de réseau \emph{primitifs}), mais il signifie
que les directions $a/b/c$ rapportées correspondent aux vecteurs de la
maille primitive telle que fournie en entrée --- qui, pour les structures
de haute symétrie (rock-salt, zinc-blende, etc.), ne coïncident pas avec les
directions cristallographiques conventionnelles ($[100]$, $[010]$, $[001]$
cubiques) auxquelles on pense intuitivement. Pour une lecture directement
comparable aux axes conventionnels, une extension naturelle consisterait à
reconventionnaliser la maille (via \texttt{pymatgen.symmetry.analyzer.
SpacegroupAnalyzer.get\_conventional\_standard\_structure}) avant
construction du graphe, ou à évaluer le poids minimal sur un ensemble plus
riche de vecteurs cibles (combinaisons simples de $a,b,c$), ce qui n'a pas
été fait dans cette première passe mais est signalé ici comme piste
d'amélioration.

Cette divergence confirme en tout cas la prémisse du projet : le descripteur
de percolation apporte une information \emph{qualitativement différente}
des agrégats ICOHP classiques, non redondante avec eux.

\section{Conclusion et perspectives}

Le pipeline complet --- requête Materials Project, calcul VASP+LOBSTER
compatible, et analyse de percolation topologiquement exacte sans
supercellule --- a été exécuté de bout en bout sur six composés réels
(3 sur le hull, 3 métastables expérimentaux), couvrant les trois grands
types de liaison (ionique, covalente, métallique). Les résultats confirment
que le descripteur de percolation capture une information distincte des
agrégats ICOHP/ICOBI usuels.

Pistes pour la suite : (i) reconventionnalisation de la maille avant
construction du graphe pour aligner les directions rapportées sur les axes
cristallographiques usuels ; (ii) extension du jeu de test à plusieurs
dizaines de composés pour une analyse statistique ; (iii) étude de
sensibilité au rayon de coupure ICOHP/ICOBI amont.

\vspace{1cm}
\noindent\textit{Code source :} \url{https://github.com/John-Akapulco/viability}

\end{document}
